% !TeX root = ../main.tex
\section{Farkas' Lemma}
This section focuses on Farkas' Lemma, an important result concerning systems of linear equations.
For notational clarity, we will follow the tradition of generally not explicitly indicating transposes.

There are many mathematically equivalent ways to state Farkas' Lemma.
We start with two versions with particularity nice geometric and algebraic interpretations.

\begin{lemma}[Farkas' Lemma, Geometric version]\label{lem:farkas-geom}
  Fix \(A \in \bF^{m \times n}\) and \(b \in \bF^m\).
  Exactly one of the following is true:
  \begin{itemize}
    \item There exists \(x \in \bF^n\), \(x \geq 0\)\footnote{\(\geq 0\) is interpreted as component wise.} such that \(Ax = b\).
    \item There exists \(y \in \bF^m\) such that \(yA \geq 0\) and \(y b < 0\).
  \end{itemize}
\end{lemma}

This says that either \(b\) lies in the positive cone generated by columns of \(A\) or there is a hyperplane separating this cone from \(b\).
\mnote{
  Thus one can prove the result above in the case of \(\bR\) using some hyperplane separation theorem for convex sets, which typically relies on topological properties of the field.
  However, Farkas' Lemma remains true for arbitrary ordered fields.
  For those cases, one needs to use an algebraic argument.
}

\begin{lemma}[Farkas' Lemma, Algebraic version]\label{lem:farkas-alg}
  Fix \(A \in \bF^{m \times n}\) and \(b \in \bF^m\).
  Exactly one of the following is true:
  \begin{itemize}
    \item There exists \(x \in \bF^n\) such that \(Ax = b\).
    \item There exists \(y \in \bF^m\) such that \(yA \geq 0\) and \(y b < 0\).
  \end{itemize}
\end{lemma}

That is, either the system \(Ax = b\) can be solved, or this system is incoherent in the sense that aggregating the equalities with weights \(y\) gives a contradiction.

\mnote{
One can show this version of Farkas' Lemma by applying the geometric version to \(\tilde{A} = [A, -A]\) and \(\tilde{b} = b\).
}

